\section{Introduction}
\label{sec:intro}

\begin{figure}[t]
  \centering
  % Placeholder for the teaser. Hand-drawn method diagrams need the comment below on the \includegraphics line.
  % \includegraphics[width=\linewidth]{figures/teaser.pdf} % fig-ok: diagram
  \fbox{\parbox{0.95\linewidth}{\centering\vspace{2.2cm}\todo{Teaser from real rendered frames of one training-scene \Look$\to$\Flip\ pair (FIG01): one row per world, \texttt{turn\_left} $\mid$ $s_0$ $\mid$ \texttt{turn\_right}, with the shared $s_0$ spanning both rows. An agent that commits at $s_0$ gives the same answer in both worlds and is wrong in one; an agent that turns first sees the kettle and answers each world correctly.}\vspace{2.2cm}}}
  \caption{\todo{Caption.} Two worlds that are indistinguishable at the start should receive the same credit for the same decision; they are allowed, and required, to diverge only after the agent has looked.}
  \label{fig:teaser}
\end{figure}

Consider an agent asked to ``select the marked object closest to the kettle'' when the kettle is out of view.
Two candidates are marked in front of it.
It could guess, or it could turn its head once and see where the kettle is.
A competent agent should spend that small observation when its current view cannot settle the question, and then act on what it sees: change its answer if the kettle turns out to be next to the other candidate, and keep it otherwise.

Vision-language models (VLMs) are increasingly trained as embodied agents with multi-turn reinforcement learning (RL) in simulation~\cite{vagen2025,glance2026,seekvln2026}.
The learning signal in these systems is computed in the one world that the episode happened to sample.
If the agent commits without looking and the kettle happens to be near its first guess, the episode succeeds and the guess is reinforced; if the kettle was elsewhere, the same decision, taken on exactly the same input, is punished.
Because the input was identical in both cases, the difference in credit carries no information the agent could have used.
In expectation the signal is still correct, but it is noisy exactly at the decision we care about, and it rewards lucky guesses along the way.
\hyp{With a finite set of training scenes the problem is sharper than noise: in our experiments, outcome-only PPO raises its training reward mainly on pairs whose answer does not depend on the hidden world, rarely looks before committing, and gains nothing in rooms it has not seen (\cref{sec:res-lucky}).}

We remove this noise at its source.
Instead of sampling one world per episode, we construct \emph{pairs} of simulated worlds that differ only in the position of one object, the anchor of the instruction.
Before the agent looks, the two worlds produce byte-identical inputs; after it looks, they may call for different answers (\Flip) or the same answer (\Stay).
We run the two worlds in lockstep: while their inputs are identical, the policy samples a single action that is executed in both, and the decision receives the same credit in both.
Two learning signals follow.
First, pre-evidence decisions are credited against the value of committing immediately, averaged over the worlds the agent cannot yet distinguish (\emph{information-set credit}).
Second, once the agent has actually seen the anchor, a consistency term asks it to switch targets in \Flip\ pairs and to keep its choice in \Stay\ pairs.

We are careful about what pairing does and does not buy.
If training worlds are sampled from the true prior, the usual single-world credit is not biased, and a same-world comparison between looking and committing, as in concurrent evidence-seeking navigation work~\cite{seekvln2026}, is also unbiased.
Our claim is narrower.
Under lockstep sampling, information-set credit is the sum over the two worlds of ``return minus the value of committing now''; it therefore combines two ingredients that can be tested separately: \emph{stratified sampling} over the hidden world, which removes the between-world component of the variance at an equal rollout budget, and a \emph{commit-now baseline} that replaces the learned critic before the evidence step.
We test them with a $2{\times}2$ design: \{single-world, paired\} sampling $\times$ \{critic, commit-now\} baseline, plus the consistency term on top.

Building the pairs requires care.
A pair is only useful if the two worlds are truly indistinguishable before the agent looks, if each has a unique and reachable target, and if the evidence actually becomes visible after one observation.
We pre-render every view the agent can reach within its observation budget in both worlds, reuse one image for both worlds only when the two renders are verified to be identical, and discard pairs that fail any validity check (\cref{sec:data}).
We also distinguish two moments that are easy to conflate: when the two worlds' inputs first differ ($\tau$), and when the agent has actually seen the anchor in its own world ($\tauseen$).
An agent that has only seen ``no kettle next to candidate 1'' has reached $\tau$ but not $\tauseen$; a commitment at that point is still a guess.

\paragraph{Contributions.}
\begin{itemize}
\item A formulation of when to look and when to change one's mind as credit assignment over information sets, with \Look/direct and \Flip/\Stay\ cases defined relative to what the agent has observed and two explicit moments, $\tau$ and $\tauseen$ (\cref{sec:method}).
\item Counterfactual world pairs in AI2-THOR~\cite{ai2thor2017} with pre-rendered view tables, verified identical inputs before the evidence step, and validity checks for physics, target uniqueness, evidence visibility, and solvability (\cref{sec:data}).
\item An analysis showing that paired credit decomposes into stratified sampling and a commit-now baseline, a controlled $2{\times}2$ design that tests each ingredient at an equal rollout budget, and a post-evidence consistency term that is applied only after the agent has seen the evidence (\cref{sec:method,sec:setup}).
\item \hyp{Experiments in unseen rooms showing that, at a matched observation budget, paired credit with post-evidence consistency reduces premature commitment and improves success on \Look\ episodes relative to outcome-only PPO and a same-world commit-now baseline, without hurting direct episodes or the unchanged task (\cref{sec:results}).} % register in CLAIMS.md before removing \hyp
\end{itemize}
